\begin{definition}
    A sequence of random varaibles $X_n$ is \textbf{stationary} if 
    \begin{equation*}
        (X_0,\ldots) \deq (X_k,\ldots)
    \end{equation*}
    for all $k\geq 0$. 
\end{definition}
\begin{remark}
    Equivalently, it means that on the product space, $X = (X_0,\ldots)$ 
    is invariant under shifts, i.e. $\theta X \deq X$. 
\end{remark}

\begin{proposition}
    If $X_n$ is stationary and $f:\R^{\Z_+}\to\R$ is measurable, then 
    $Y_n = f(X_n,X_{n+1},\ldots)$ is stationary. 
\end{proposition}
\begin{proof}
    It suffices to show that for all $n\geq 0$, $(Y_0,\ldots,Y_n) \deq (Y_{k},\ldots,Y_{k+n})$. 
    Indeed, for all measurable $B\subset\R^{n+1}$, 
    \begin{equation*}
        \P((Y_0,\ldots,Y_n)\in B) = \P((X_0,\ldots)\in A) = \P((X_k,\ldots)\in A) = \P((Y_k,\ldots,Y_{k+n})\in B) 
    \end{equation*}
    for some measurable $A\subset\R^{\Z_+}$. 
\end{proof}

\begin{definition}
    Let $(\Omega,\F,\P)$ be a probability space. A measurable mapping 
    $T:\Omega\to \Omega$ is called \textbf{measure-preserving} if 
    $\P(T^{-1}(A)) = \P(A)$ for all $A\in\F$. 
\end{definition}
\begin{remark}
    By the approximation argument, this also implies that 
    $\E\sbrc{X} = \E\sbrc{XT}$ for all measurable $X\geq 0$ 
    or $\E\sbrc{\abs{X}}<\infty$.    
\end{remark}

\begin{proposition}
    Let $\Omega$ be a probability space, $X$ be a random variable on 
    $\Omega$ and $T:\Omega\to \Omega$ be measurable. Then if $T$ is 
    measure-preserving, $(XT^n)$ is stationary. Conversely, if $\F = \sigma(X)$ 
    and $XT^n$ is stationary, then $T$ is measure-preserving.
\end{proposition}
\begin{proof}
    If $T$ is measure-preserving, then $\P(XT^n\in B) = \P(T^n\omega \in X^{-1}(B)) 
    = \P(\omega\in X^{-1}(B)) = \P(X\in B)$. Hence $XT^n$ is stationary. 
    Conversely, for every $A\in\F$, there is $B$ such that $A = X^{-1}(B)$. 
    Since $(XT^n)$ is stationary, $\P(T^{-1}(A)) = \P(T^{-1}X^{-1}(B)) 
    = \P(XT\in B) = \P(X\in B) = \P(A)$. Hence $T$ is measure-preserving. 
\end{proof}

\begin{proposition}\label{prop:measure_preserving_exist}
    Let $X_n$ be stationary in $S$. Consider $\Omega = S^{\Z_+}$. 
    Then the shift operator on $\Omega$ is measure-preserving.  
\end{proposition}
\begin{proof}
    For $\omega = (\omega_0,\ldots)$, we can rewrite $X_n$ on $\Omega$ 
    with $X_n(\omega) = \omega_n$ and $X_n = \omega_{n} = X_0\theta^n$ by 
    definition. Since $X_n$ is stationary, $\theta$ is measure-preserving 
    on $\Omega$. 
\end{proof}

\begin{definition}
    Suppose that $T:\Omega\to\Omega$ is measure-preserving. A set $A\in\F$ 
    is \textbf{invariant} if $T^{-1}A = A$ almost surely, i.e., 
    $\P((T^{-1}A)\triangle A) = 0$. The collection of invariant 
    sets $\I$ is called the \textbf{invariant $\sigma$-algebra}. 
\end{definition}
\begin{remark}
    $\I$ is indeed a $\sigma$-algebra. First it is clear that $\varnothing\in\F$. 
    If $A\in\I$, then 
    \begin{equation*}
        \begin{split}
            \P((T^{-1}A^c)\triangle A^c) &= \P((T^{-1}A^c - A^c)\cup (A^c - T^{-1}A^c)) \\
            &= \P((A-T^{-1}A)\cup (T^{-1}A - A)) = \P((T^{-1}A)\triangle A) = 0. 
        \end{split}
    \end{equation*} 
    Hence $A^c\in\I$. Finally, if $A = \cup_n A_n$ with $A_n\in\F$, 
    $T^{-1}A = \cup_n(T^{-1}A_n) = \cup_nA_n = A$ almost surely. 
    Hence $\I$ is a $\sigma$-algebra. 
\end{remark}

\begin{definition}
    Let $T:\Omega\to\Omega$ be measure-preserving with invariant $\sigma$-algebra 
    $\I$. $T$ is said to be \textbf{ergodic} if $\I$ is trivial. A 
    stationary sequnece $X_n$ is \textbf{ergodic} if the associated 
    measure-preserving mapping is ergodic.
\end{definition}

\begin{proposition}
    Let $X_n$ be independent and identically distributed. Then $X_n$ 
    is ergodic. 
\end{proposition}
\begin{proof}
    Notice that the invariant $\sigma$-algebra $\I$ generated by the shift 
    operator satisfies that $\I\subset\T$, where $\T$ is the tail 
    $\sigma$-algebra. Indeed, if $A\in\I$, then $(\theta^n)^{-1}A\theta^{-1}A = A$ 
    almost surely. This implies that $A\in\sigma(X_n,\ldots)$ for all $n$. 
    Hence $A\in\T$ and $\I\subset\T$. Now by the Kolmogorov's zero 
    one law, $\T$ is trivial and so is $\I$. 
\end{proof}

\begin{theorem}
    Let $X_n$ be an irreducible Markov chain on a countable space starting with the stationary 
    distribution $\pi$. Then $X_n$ is ergodic. 
\end{theorem}
\begin{proof}
    Suppose that $A\in\I$ and $\F_n = \sigma(X_0,\ldots,X_n)$. Then 
    $\one_A\theta = \one_A$. Now 
    \begin{equation*}
        \E_\pi\sbrc{\one_A|\F_n} = \E_\pi\sbrc{\one_A\theta^n|\F_n} 
        = \E_{X_n}\sbrc{\one_A}.  
    \end{equation*}
    Denote $h(x) = \E_x\sbrc{\one_A}$. By Levy's zero one law, 
    $\E_\pi\sbrc{\one_A|\F_n}\to\one_A$ almost surely. Note that 
    $h(X_n)$ is a bounded martingale. By the martingale convergence 
    theorem, $h(X_n)$ converges almost surely to some random 
    variable. If $h$ is not a constant function, there are $x,y$ 
    such that $h(x)\neq h(y)$ and $X_n$ visits $x$ and $y$ infinitely 
    many times by the positive recurrence, contradicting to the fact 
    that $h(X_n)$ converges. Thus $h$ is constant. Hence $h = 0$ or $1$. 
    This implies that $\P(A|\F_n) = 0$ or $1$ and $\P(A) = 0$ or $1$. 
    We conclude that $\I$ is trivial.  
\end{proof}

\begin{example}
    Fix $\theta\in [0,1]$. Consider a dynamical system 
    starting from $\omega_0\in[0,1]$ and $\omega_{n+1} = \omega_n 
    + \theta\mod 1$. Suppose that the measure is the usual 
    Lebesgue measure. Then the process is ergodic if and only if 
    $\theta\notin\Q$. 
    
    First suppose that $\theta\in\Q$, then $\omega_n$ returns to $\omega_0$ 
    in finite number of time, say $\omega_n = \omega_0$. Then we may 
    find a tiny interval $I_i$ centered at $\omega_i$ for each $i\leq n$.  
    The intervals can be selected such that $\P(\cup_{i\leq n}I_i)\in (0,1)$. 
    But $\cup_{i\leq n}I_i$ is invariant. Hence the process is not 
    ergodic. 
    
    Next, we show that if $\theta\notin\Q$, then the process is ergodic. 
    For any $A\subset [0,1]$, if $A$ is invariant, then $\one_A(\omega) 
    = \one_A(\omega + \theta)$. By the Fourier expansion, 
    \begin{equation*}
        \sum_{n\in\Z} c_n e^{2\pi i\omega n} = \one_A(\omega) = \one_A(\omega + \theta) 
        = \sum_{n\in\Z} c_n e^{2\pi in(\omega + \theta)} = \sum_{n\in\Z} \pth{c_n e^{2in\theta\pi}} e^{2\pi in\omega}
    \end{equation*}
    where 
    \begin{equation*}
        c_n = \int_0^{1} \one_A(\omega) e^{-2\pi in\omega}d\omega. 
    \end{equation*}
    By the uniqueness of Fourier expansion, since $\theta$ is irrational, 
    $c_n = 0$ for all $n\neq 0$. Thus $c_0 = \one_A$ and $\P(A) = 0$ 
    or $1$. 
\end{example}

\begin{example}
    Let $\Omega = [0,1]$ and $T:\omega\mapsto\set{2\omega}$ where $\set{x}$ 
    denote the fractional part of $x$. Then $T$ is ergodic with respect 
    to the Lebesgue measure. To see this, note that for all $\omega\in[0,1]$, 
    $\sum_{i=1}^\infty2^{-i}\omega_i$ where $\omega_i\iidsim\Ber(1/2)$ 
    follows the Lebesgue measure. $T:(\omega_1,\ldots) = (\omega_2,\ldots)$ 
    is a shift of independent and identically distributed random variables 
    and hence is ergodic. 
\end{example}

\begin{lemma}[Hopf's Maximal Ergodic Lemma]
    Let $T$ be a measure-preserving map and $X\in\L^1$. Suppose $X_k(\omega) 
    = X(T^k\omega)$. Let $S_k = \sum_{j=0}^{k-1}X_j$ and $M_n = 
    \max\set{0,S_1\ldots,S_n}$. Then 
    \begin{equation*}
        \E\sbrc{X\one\set{M_n>0}} \geq 0. 
    \end{equation*}
\end{lemma}
\begin{proof}
    Write $M_k(T\omega) = \max\set{0,X_1,\ldots,X_1+\cdots+X_k}$. 
    Then 
    \begin{equation*}
        X_0 + M_k(T\omega) = \max\set{X_0,\ldots,X_0+X_1+\cdots+X_k}\geq S_j
    \end{equation*} 
    for all $j = 1,\ldots k+1$. This implies that 
    $\sbrc{X_0 + M_k(T\omega)}\one\set{M_k>0}\geq M_k\one\set{M_k>0}$ 
    and $X_0\one\set{M_k>0}\geq (M_k - M_kT)\one\set{M_k>0}$. Notice that 
    by the stationarity, 
    \begin{equation*}
        0 = \E\sbrc{M_k - M_kT} = \E\sbrc{(M_k - M_kT)\one\set{M_k=0}} 
        + \E\sbrc{(M_k - M_kT)\one\set{M_k>0}}. 
    \end{equation*}
    We know that $\E\sbrc{(M_k - M_kT)\one\set{M_k=0}} \leq 0$ and hence 
    $\E\sbrc{(M_k - M_kT)\one\set{M_k>0}}\geq 0$. Thus 
    \begin{equation*}
        \E\sbrc{X\one\set{M_n>0}}\geq \E\sbrc{(M_k - M_kT)\one\set{M_k>0}}\geq 0.
    \end{equation*}
    The proof is complete. 
\end{proof}

\begin{theorem}[Birkhoff's Pointwise Ergodic Theorem]
    Let $X\in\L^1$ and $T$ be a measure-preserving mapping with the invariant 
    $\sigma$-algebra $\I$. Then 
    \begin{equation*}
        \frac{1}{n}\sum_{k=0}^{n-1}X(T^k\omega)\to \E\sbrc{X|\I}
    \end{equation*}
    almost surely and in $\L^1$. If $T$ is ergodic, then $\E\sbrc{X|\I} 
    = \E\sbrc{X}$. 
\end{theorem}
\begin{proof}
    Since $\E\sbrc{X|\I}$ is $\I$-measurable, $X(T^k\omega) - \E\sbrc{X|\I}$ 
    is again stationary and hence we may without loss of generality 
    assume that $\E\sbrc{X|\I}=0$. Put 
    \begin{equation*}
        \bar{X} = \limsup_{n\to\infty}\frac{1}{n}\sum_{k=0}^{n-1} X_k
    \end{equation*}
    where $X_k = XT^k$. Observe that 
    \begin{equation*}
        \bar{X}(T\omega) = \limsup_{n\to\infty}\frac{1}{n}\sum_{k=1}^{n} X_k 
        = \limsup_{n\to\infty}\frac{n+1}{n}\cdot\frac{1}{n+1}\sum_{k=0}^{n} X_k  - \frac{1}{n}X_0 
        = \bar{X}.
    \end{equation*}
    Thus $\bar{X}$ is $\I$-measurable. We claim that $\P(\bar{X}>\epsilon) = 0$ 
    for all $\epsilon>0$. Let $D = \set{\bar{X}>\epsilon}\in\I$ and 
    $X^* = (X-\epsilon)\one_D$. Set $S_n = X^* + X^*T + \cdots + X^*T^{n-1}$ 
    and $M_n = \max\set{0,S_1,\ldots,S_n}$. Write $F_n = \set{M_n>0}$ 
    and $F_n\nearrow F \coloneq \set{\sup_{k}k^{-1}S_k>0}$. Since $D$ is 
    invariant, $X^*T^n = (X_n-\epsilon)\one_D$. By the definition 
    of $D$, $F = D\cap \set{\sup_{k}k^{-1}S_k>0} = D$. By Hopf's 
    maximal ergodic lemma, $\E\sbrc{X^*\one_{F_n}}\geq 0$. Since 
    $\abs{X^*}\leq \abs{X}+\epsilon\in\L^1$, $\E\sbrc{X^*\one_{F_n}}
    \to\E\sbrc{X^*\one_F}$ by LDCT. Thus $\E\sbrc{X^*\one_F}\geq 0$. 
    Now 
    \begin{equation*}
        \E\sbrc{X\one_D} - \epsilon\P(D) = \E\sbrc{(X-\epsilon)\one_D} = \E\sbrc{X^*\one_D}\geq 0. 
    \end{equation*}
    But $\E\sbrc{X\one_D} = \E\sbrc{\E\sbrc{X\one_D|\I}} = \E\sbrc{\one_D\E\sbrc{X|\I}} = 0$. 
    Hence $\P(D) = 0$. The claim follows and thus $\limsup_{n\to\infty}\frac{1}{n}\sum_{k=0}^{n-1} X_k\leq 0$ 
    almost surely. 

    Similarly, one can show that $\liminf_{n\to\infty}\frac{1}{n}\sum_{k=0}^{n-1} X_k\geq 0$ 
    almost surely. The almost sure convergence follows. To prove the 
    $\L^1$ convergence, write $X = Y + Z$ where $Y = X\one\set{\abs{X}\leq M}$ 
    and $Z = X\one\set{\abs{X} > M}$. Taking $n\to\infty$ and then $M\to\infty$, 
    \begin{equation*}
        \frac{1}{n}\sum_{k=0}^{n-1}YT^k\to \E\sbrc{X\one\set{\abs{X}\leq M}|\I}\to\E\sbrc{X|\I}
    \end{equation*}
    in $\L^1$ by the bounded convergence theorem and LMCT. Also, 
    \begin{equation*}
        \E\sbrc{\frac{1}{n}\sum_{k=0}^{n-1}\abs{ZT^k}} = \E\sbrc{\abs{Z}} = \E\sbrc{X\one\set{\abs{X}>M}}\to 0
    \end{equation*}  
    as $M\to\infty$ by LDCT. This completes the proof. 
\end{proof}

\begin{definition}
    Let $\H$ be a Hilbert space. A linear operator $S$ on $\H$ is \textbf{isometric} 
    if $\inp{Sx}{Sy} = \inp{x}{y}$ for all $x,y\in\H$. 
\end{definition}

\begin{theorem}[Von Neumann's Mean Ergodic Theorem]
    Let $\H$ be a Hilbert space, $S:\H\to\H$ be an isometric linear 
    operator and $P$ be the orthogonal projection onto $\ker(S-I)$. Then 
    \begin{equation*}
        \frac{1}{n}\sum_{k=0}^{n-1}S^kx\to Px
    \end{equation*}
    in the norm on $\H$ for all $x\in\H$. 
\end{theorem}
\begin{proof}
    We first claim that $\H = \ker(S-I)\oplus \overline{R(S-I)}$. To 
    prove the claim, it suffices to show that $R(S-I)^\perp = \overline{R(S-I)}^\perp = \ker(S-I)$. 
    Indeed, if $x\in \ker(S-I)$, then $Sx = x$ and for all $y\in R(S-I)$,  
    $y = Sz - z$, 
    \begin{equation*}
        \inp{x}{y} = \inp{x}{Sz-z} = \inp{x}{Sz} - \inp{x}{z} = \inp{Sx}{z} - \inp{x}{z} = 0. 
    \end{equation*}
    Hence $x\in R(S-I)^\perp$. Conversely, if $x\in R(S-I)^\perp$, 
    \begin{equation*}
        0 = \inp{x}{Sx-x} = \inp{\frac{1}{2}(x-Sx) + \frac{1}{2}(Sx + x)}{Sx-x} 
        = -\frac{1}{2}\norm{Sx-x}_\H^2 + \frac{1}{2}\pth{\norm{Sx}_\H^2 - \norm{x}_\H^2}.
    \end{equation*}
    Since $S$ is isometric, $\norm{Sx}_\H = \norm{x}_\H$ and $Sx = x$. Thus 
    $x\in\ker(S-I)$. We see that $R(S-I)^\perp = \ker(S-I)$. 

    Now let $A_n = \frac{1}{n}\sum_{k=0}^{n-1}S^k$. For all $x\in\H$, write $x = y+z$ where 
    $y = Px\in\ker(S-I)$ and $z\in\overline{R(S-I)}$. $A_ny = \frac{1}{n}\sum_{k=0}^{n-1}S^ky = y$ 
    since $Sy = y$. Also, there is $v$ such that $z = Sv-v$.
    \begin{equation*}
        A_nz = \frac{1}{n}\sum_{k=0}^{n-1}S^k(Sv-v) = \frac{1}{n}(S^nv - v)\to 0
    \end{equation*}
    since $\norm{S^n-I}\leq \norm{S^n}+1 = 2$. We conclude that $A_nx\to x$ 
    as $n\to\infty$. 
\end{proof}

\begin{corollary}
    Let $X\in\L^2(\Omega,\P)$, $T:\Omega\to\Omega$ be measure-preserving 
    and $X_k(\omega) = X(T^k\omega)$. If $\I$ is the invariant $\sigma$-algebra, 
    then 
    \begin{equation*}
        \frac{1}{n}\sum_{k=0}^{n-1}X_k\to \E\sbrc{X|\I}
    \end{equation*}
    in $\L^2$.
\end{corollary}
\begin{proof}
    Define $S:X\mapsto XT$ to be a linear operator. Note that 
    \begin{equation*}
        \int X(T\omega)Y(T\omega)d\P = \int XYd(\P\circ T^{-1}) = \int XYd\P.
    \end{equation*}
    Hence $S$ is an isometry. Notice that $\ker(S-I) = \L^2(\I)$. 
    By the von Neumann's mean ergodic theorem, $\frac{1}{n}\sum_{k=0}^{n-1}X_k\to PX$ 
    where $P$ is the projection onto $\L^2(\I)$. By definition, 
    $PX = \E\sbrc{X|\I}$. The corollary follows. 
\end{proof}

\begin{definition}
    A measure-preserving transformation $T$ on $(\Omega,\F,\P)$ 
    is called \textbf{mixing} if $\lim_{n\to\infty}\P(A\cap T^{-n}B) 
    = \P(A)\P(B)$ for all $A,B\in\F$. 
\end{definition}

\begin{definition}
    A measure-preserving transformation $T$ on $(\Omega,\F,\P)$ 
    is called \textbf{average mixing} if 
    \begin{equation*}
        \lim_{n\to\infty}\frac{1}{n}\sum_{k=0}^{n-1}\P(A\cap T^{-n}B) 
        = \P(A)\P(B). 
    \end{equation*} 
\end{definition}
\begin{remark}
    Clearly mixing implies average mixing. This follows 
    from the fact that a sequence converges, then it 
    converges in Cesaro's sense.  
\end{remark}

\begin{proposition}
    Suppose that $T$ is measure-preserving on $(\Omega,\F,\P)$. Show that 
    $T$ is ergodic if and only if $T$ is average mixing. 
\end{proposition}
\begin{proof}
    Suppose that $T$ is ergodic. By the Birkhoff's ergodic theorem, 
    \begin{equation*}
        \frac{1}{n}\sum_{k=0}^{n-1}\one_{T^{-1}B}
        = \frac{1}{n}\sum_{k=0}^{n-1}\one_B(T\omega) \to 
        \E\sbrc{\one_B|\I} = \P(B)
    \end{equation*}
    almost surely. Hence 
    \begin{equation*}
        \frac{1}{n}\sum_{k=0}^{n-1}\one_{A\cap T^{-1}B} 
        = \one_A\sbrc{\frac{1}{n}\sum_{k=0}^{n-1}\one_{T^{-1}B}} 
        \to \one_A\P(B)
    \end{equation*}
    almost surely. Clearly the left hand side is bounded, the bounded convergence 
    theorem implies that 
    \begin{equation*}
        \frac{1}{n}\sum_{k=0}^{n-1}\P(A\cap T^{-1}B)\to\E\sbrc{\one_A\P(B)} = \P(A)\P(B).
    \end{equation*}

    Conversely, let $A\in\I$. 
    \begin{equation*}
        \frac{1}{n}\sum_{k=0}^{n-1}\P(A\cap T^{-k}B) 
        = \frac{1}{n}\sum_{k=0}^{n-1}\P(T^{-k}A\cap T^{-k}B)  
        = \frac{1}{n}\sum_{k=0}^{n-1}\P(T^{-k}(A\cap B))  
        = \P(A\cap B). 
    \end{equation*}
    On the other hand, $\frac{1}{n}\sum_{k=0}^{n-1}\P(A\cap T^{-k}B)\to\P(A)\P(B)$ 
    almost surely. Hence $\P(A\cap B) = \P(A)\P(B)$. This means that $\I$ is 
    independent of $\F$ and in particular $\P(A) = \P(A)^2$ implies that $\P(A) = 0$ 
    or $1$. Thus $\I$ is trivial and $T$ is ergodic. 
\end{proof}
\begin{remark}
    Consider $T:\omega\to\set{\omega+\theta}$ with $\theta\notin\Q$. 
    $\set{\cdot}$ denote the fractional part. We have seen that $T$ 
    is ergodic. We claim that it is not mixing. 

    Consider $A = B = [0,1/2]$. Observe that there is a subsequence $n_k$ such 
    that $\set{n_k\theta}\to 0$ as $k\to\infty$. To see this, fix $\epsilon>0$. 
    Let $N\in\N$ with $1/N<\epsilon$. By Dirichlet's theorem, there are integers 
    $p,q$, $1\leq q\leq N$ such that $\abs{q\theta - p}<1/N<\epsilon$. Hence 
    $\set{q\theta}$ satisfies either $\set{q\theta}<\epsilon$ or $\set{q\theta}>1-\epsilon$. 
    In the later case, $\set{mq\theta} = 1-m(1-\set{q\theta})$ for $m$ such that 
    $1-m(1-\set{q\theta})>0$ and we may find $m$ such that $1-m(1-\set{q\theta})<\epsilon$ 
    since $1-\set{q\theta}<\epsilon$. We conclude that we can always find $r\in\N$ 
    such that $\set{r\theta}<\epsilon$. Now pick $\epsilon_k\searrow 0$, we can find 
    associate subsequence $n_k$ such that $\set{n_k\theta}\to 0$. For such subsequence, 
    \begin{equation*}
        \P(A\cap T^{-n_k}B) = \frac{1}{2}-\set{n_k\theta}\to \frac{1}{2}\neq \frac{1}{4} = \P(A)\P(B). 
    \end{equation*}
    Thus $T$ is not mixing. 
\end{remark}

\begin{lemma}[Fekete]
    Let $a_n\in\R$ be such that $a_k+a_l\geq a_{k+l}$. Then 
    \begin{equation*}
        \lim_{n\to\infty} \frac{a_n}{n} = \inf_{n}\frac{a_n}{n}. 
    \end{equation*}
\end{lemma}
\begin{proof}
    Since $\liminf_{n\to\infty}\frac{a_n}{n}\geq \inf_{n}\frac{a_n}{n}$. 
    It suffices to show that $\limsup_{n\to\infty}\frac{a_n}{n}\leq 
    \inf_{n}\frac{a_n}{n}$. Let $n = km+l$ with $0\leq l\leq m-1$. 
    Then $a_{km}+a_l\geq a_n$ and $ka_m\geq a_{km}$. Thus $ka_m + a_l\geq a_n$. 
    Dividing by $n$, 
    \begin{equation*}
        \frac{km}{n}\frac{a_m}{m} + \frac{a_l}{n}\geq \frac{a_n}{n}. 
    \end{equation*} 
    Fix $m,l$ and let $n,k\to\infty$. Then 
    \begin{equation*}
        \frac{a_m}{m}\geq \limsup_{n\to\infty}\frac{a_n}{n}
    \end{equation*}
    for all $m$. Thus $\inf_{m}\frac{a_m}{m}\geq \limsup_{n\to\infty}\frac{a_n}{n}$. 
    The lemma follows. 
\end{proof}

\begin{theorem}[Subadditive Ergodic Theorem, Liggett]
    Suppose that $X_{m,n}$ with $0\leq m<n$ satisfies 
    \begin{thmenum}
        \item $X_{0,m}+X_{m,n}\geq X_{0,n}$ for $0\leq m <n$. 
        \item Given $k\geq 1$, $\set{X_{nk,(n+1)k}}$ is stationary. 
        \item The distribution of $X_{m,m+k}$ does not depend on $m$. 
        \item $\E\sbrc{X_{0,1}^+}<\infty$ and $\E\sbrc{X_{0,n}}\geq -cn$ 
        for some constant $c$. 
    \end{thmenum}
    Then 
    \begin{thmenum}
        \item $\gamma\coloneq\lim_{n\to\infty}\frac{\E\sbrc{X_{0,n}}}{n} = \inf_{n}\frac{\E\sbrc{X_{0,n}}}{n}>-\infty$. 
        \item $\frac{X_{0,n}}{n}\to X$ almost surely and in $\L^1$ for some $X\in\L^1$ with $\E\sbrc{X}=\gamma$.
        \item If the sequence in condition (b) is ergodic, then $X = \gamma$. 
    \end{thmenum}
\end{theorem}
\begin{proof}
    For (a), notice that by conditions (a), (c), 
    $\E\sbrc{X_{0,m}} + \E\sbrc{X_{0,n-m}}\geq \E\sbrc{X_{0,n}}$. 
    Applying Fekete's lemma shows that 
    \begin{equation*}
        \gamma = \lim_{n\to\infty}\frac{\E\sbrc{X_{0,n}}}{n} = \inf_{n}\frac{\E\sbrc{X_{0,n}}}{n}
    \end{equation*}
    Furthermore, by condition (d), $\gamma\geq -c>-\infty$ for some constant $c$. 

    Next, define 
    \begin{equation*}
        \bar{X} = \limsup_{n\to\infty}\frac{X_{0,n}}{n}, \quad 
        \underbar{X} = \liminf_{n\to\infty}\frac{X_{0,n}}{n}. 
    \end{equation*}
    We claim that $\E\sbrc{\bar{X}}\leq \gamma$ 
    and if the stationary process in condition (b) 
    is ergodic, $\bar{X}\leq \gamma$ almost surely. 
    To prove this, fix $k\geq 1$. By the condition (a), 
    \begin{equation*}
        X_{0,kn}\leq \sum_{j=1}^nX_{k(j-1),kj}. 
    \end{equation*} 
    Using the Birkhoff's ergodic theorem, 
    \begin{equation*}
        \frac{1}{n}\sum_{j=1}^nX_{k(j-1),kj} 
        \to X_k 
    \end{equation*}
    almost surely and in $\L^1$ for some 
    random variable $X_k$. Hence
    \begin{equation*}
        \E\sbrc{\limsup_{n\to\infty}\frac{X_{k,n}}{kn}}\leq \frac{\E\sbrc{X_k}}{k}. 
    \end{equation*}
    Notice also that $X_{kn+j}\leq X_{kn} + X_{kn,kn+j}$. 
    By condition (c), the distribution of $X_{kn,kn+j}$ only depends 
    on $j$. By condition (d), $X_{kn,kn+j}$ has finite 
    first moment. By the Borel-Cantelli lemma, one can show that 
    \begin{equation*}
        \lim_{n\to\infty}\frac{X_{kn,kn+j}}{n} = 0
    \end{equation*}
    almost surely for all $j$. Thus $\E\sbrc{\bar{X}}
    \leq \E\sbrc{X_k}/k\to\gamma$. If the process 
    is ergodic, then $X_k$ is deterministic and 
    $\bar{X}\leq \gamma$ almost surely. 

    Now we shall claim that $\E\sbrc{\underbar{X}}\geq \gamma$. 
    Let $U_n$ be independent of $X_{k,n}$ and distributed 
    uniformly on $\set{1,\ldots,n}$. Set $Y^n_k = X_{0,k+U_n}-X_{0,k+U_n-1}$. 
    Notice that 
    \begin{equation*}
        \E\sbrc{Y^n_k} = \frac{1}{n}\sum_{j=1}^n\E\sbrc{X_{0,k+j}-X_{0,k+j-1}} 
        = \frac{1}{n}\pth{\E\sbrc{X_{0,k+n}} - \E\sbrc{X_{0,k}}}
    \end{equation*}
    and 
    \begin{equation*}
        \E\sbrc{(Y^n_k)^+} = \frac{1}{n}\sum_{j=1}^n\E\sbrc{(X_{0,k+j}-X_{0,k+j-1})^+} 
        \leq  \frac{1}{n}\sum_{j=1}^n\E\sbrc{X_{k+j-1,k+j}^+} = \E\sbrc{X_{0,1}^+}. 
    \end{equation*}
    Hence $\sup_{n}\E\sbrc{\abs{Y_k^n}}<\infty$ and 
    $\lim_{n\to\infty}\E\sbrc{Y_k^n} = \gamma$ for all $k\geq 1$. Now 
    there is a subsequence, along which the joint distribution $\set{Y_k^{n_i}}$ 
    converges to the joint distribution, say $\set{Y_k}$. For any $f\in C_b(\R^\N)$ 
    depending only on finitely many coordinates, we have 
    \begin{equation*}
        \E\sbrc{f(Y_1,\ldots)} 
        = \lim_{i\to\infty} \frac{1}{n_i}\sum_{j=1}^{n_i}\E\sbrc{f(X_{0,j+1}-X_{0,j}, X_{0,j+2} - X_{0,j+1},\ldots)}. 
    \end{equation*}  
    Thus $\set{Y_k}$ is a stationary sequence. Using condition (a), 
    \begin{equation*}
        Y^n_1 = X_{0,U_n+1} - X_{0,U_n} \leq X_{U_n,U_n+1} \deq X_1. 
    \end{equation*}
    Hence by condition (d), $\set{(Y^n_1)^+}$ is uniformly integrable. 
    By Fatou's lemma, 
    \begin{equation*}
        \E\sbrc{Y_1}\geq \limsup_{n\to\infty}\E\sbrc{Y_1^n} = \gamma. 
    \end{equation*}
    Using Birkhoff's ergodic theorem again shows that 
    \begin{equation*}
        \frac{1}{n}\sum_{k=1}^nY_k \to Y
    \end{equation*}
    almost surely for some $Y$ with $\E\sbrc{Y} = \E\sbrc{Y_1}\geq \gamma$. 
    To establish the claim, it remains to show that $\underbar{X}$ is 
    stochastically larger than $Y$.\footnote{A real-valued random variable 
    $X$ is said to be stochastically larger than $Y$ if $\P(X>a)\geq \P(Y>a)$ 
    for all $a\in\R$.} To see this, note that for every increasing $f\in C_b(\R^\N)$ 
    depending only on finitely many coordinates, 
    \begin{equation*}
        \begin{split}
            \E\sbrc{f(Y_1, Y_1+Y_2,Y_1+Y_2+Y_3,\ldots)} 
            &= \lim_{i\to\infty}\frac{1}{n_i}\sum_{j=1}^{n_i}\E\sbrc{f(X_{0,j+1}-X_{0,j}, X_{0,j+2} - X_{0,j+1},\ldots)} \\ 
            &\leq \lim_{i\to\infty}\frac{1}{n_i}\sum_{j=1}^{n_i}\E\sbrc{f(X_{j,j+1}, X_{j+1,j+2},\ldots)} 
            = \E\sbrc{f(X_{0,1},\ldots)}. 
        \end{split}
    \end{equation*}
    Hence the sequence $X_1,\ldots$ is stochastically larger 
    than $Y_1,Y_1+Y_2,\ldots$. It follows that $\underbar{X}$ 
    is stochastically larger than $Y$. This proves the almost sure 
    convergence in (b) and (c).

    Finally, for the $\L^1$ convergence, let $X = \bar{X} = \underbar{X}$. 
    $\E\sbrc{X} = \gamma$. Since 
    \begin{equation*}
        X_{0,n}\leq \sum_{j=1}^nX_{(j-1),j}, 
    \end{equation*}
    $\frac{1}{n}X_{0,n}^+$ is uniformly integrable. Thus 
    $\E\sbrc{(n^{-1}X_{0,n}-X)}^+\to 0$. Now 
    \begin{equation*}
        0 = \lim_{n\to\infty}\E\sbrc{\frac{1}{n}X_{0,n}-X}\leq \lim_{n\to\infty}\E\sbrc{\pth{\frac{1}{n}X_{0,n}-X}^+} - \E\sbrc{\pth{\frac{1}{n}X_{0,n}-X}^-}.
    \end{equation*}
    Thus the negative part also converges to $0$. This gives the 
    desired $\L^1$ convergence. 
\end{proof}